\documentclass[10pt,a4paper]{amsart}
\usepackage[margin=19mm]{geometry}
\usepackage{amsmath,amssymb,mathtools}
\usepackage[scr=rsfs,cal=euler]{mathalfa}
\usepackage{xfrac}
\usepackage[unicode,hidelinks,bookmarks=false]{hyperref}
\newcommand{\ie}{i.e.\ }
\newcommand{\cf}{cf.\ }
\newcommand{\resp}{respectively\ }
\newcommand{\Subsection}{\subsection}
\providecommand{\nobreakdash}{\nobreak\hbox{-}}
\setlength{\emergencystretch}{2em}
\multlinegap0pt
\usepackage{amssymb}
\usepackage{float}
\usepackage[shortlabels]{enumitem}
\newtheorem{theorem}{Theorem}
\title{Peakon solutions and analytical properties for the Camassa-Holm type equations with quadratic nonlinearities}
\author{Yonghong Chen, Zhijun Qiao, Mingxuan Zhu$^\ast$}
\date{Source: arXiv:2605.20773v1 (CC BY 4.0)}
\begin{document}
\maketitle

\noindent\emph{Source and licence.} Peakon solutions and analytical properties for the Camassa-Holm type equations with quadratic nonlinearities. Version arXiv:2605.20773v1 (CC BY 4.0), distributed by arXiv under the Creative Commons Attribution 4.0 International licence, http://creativecommons.org/licenses/by/4.0/, as recorded in the arXiv licence metadata for this exact version and shown on the abstract page https://arxiv.org/abs/2605.20773v1. Full licence notice: \texttt{licenses/arxiv-2605.20773-CC-BY-4.0.txt}.\par\smallskip

\noindent\textit{Editorial scope.} Counted unit: the article's theorem bls1 (source0.tex lines 600--606). The article's complete original proof of it is delivered with it (source0.tex lines 607--660, one \texttt{proof} environment of 54 lines). No mathematics is written, repaired or re-derived: the statement and the proof are the article's own lines, in the article's own notation, and no retained line is substituted or re-flowed.  Retained uncounted as the source-local context the proof needs: lines 534--541; lines 586--588; lines 593--596.  Nothing was omitted from any retained span, so the package carries no omission record.  No retained line contains a citation, so the wrapper carries no bibliography and no bibliography entry.  No retained line contains a figure environment, an image-inclusion command or drawing code, so no image is delivered and the package declares no asset dependency.  Editorial formatting only, in addition: an AMS \texttt{amsart} preamble that restates the article's own declarations and macros used by the retained lines, namely the environments \texttt{theorem} on separate counters, together with the standard packages those lines need.  The retained environments print in this document as Theorem 1 in order of appearance, whereas the article prints the same items with the numbering of its own full document.  The complete native source of the article as distributed in the arXiv e-print for this exact version is bundled unchanged as \texttt{sources/201017/source0.tex}.
\begin{equation}\label{Cauchyp}
\left\{
    \begin{aligned}
        &u_t + \lambda_1 u^2 + \lambda_2 u u_x + G\ast(\lambda_3 u^2 + \lambda_4 u_x^2) + \partial_x \left(  G \ast (\lambda_5 u^2 + \lambda_6 u_x^2) \right) = 0, \quad x\in\mathbb{R},\ t>0, \\
        &u(x,0)=u_0, \quad x\in\mathbb{R}.
    \end{aligned}
\right.
\end{equation}

\begin{equation}\label{KP1}
\|fg\|_{H^r}\leq c(\|f\|_{L^\infty}\|g\|_{H^r}+\|f\|_{H^r}\|g\|_{L^\infty}),
\end{equation}

\begin{equation}\label{KP2}
\|[\Lambda^s,f]g\|_{L^2}\leq c(\|\partial_xf\|_{L^\infty}\|\Lambda^{r-1}g\|_{L^2}
+\|\Lambda^sf\|_{L^2}\|g\|_{L^\infty}),
\end{equation}

\begin{theorem}\label{bls1}
Let $u_0(x)\in H^s(\mathbb{R})$, $s>\frac{3}{2}$ and let $T$ be the existence time of the solution $u(x,t)$ to  the Cauchy problem \eqref{Cauchyp} with the initial data $u_0(x)$. Suppose that there exist $M > 0$ such that
\begin{eqnarray}\label{blsc}
\|u(\cdot,t)\|_{L^\infty}+\|u_x(\cdot,t)\|_{L^\infty}\leq M, \quad t\in [0,T),
\end{eqnarray}
then the $H^s$-norm of $u(x,t)$ does not blow up on $[0,T)$.
\end{theorem}

\begin{proof}
Let $\Lambda=(1-\partial_x^2)^{\tfrac{1}{2}}$. Applying $\Lambda^s$ on \eqref{Cauchyp}, multiplying by
$\Lambda^s u$ and integrating with respect to $x$ over $\mathbb{R}$, we have
\begin{equation}\label{Hsblow-up}
\begin{aligned}
\frac{1}{2}\frac{\mathrm{d}}{\mathrm{d}t}\|u\|_{H^s}^2
&=-\lambda_1\int_{\mathbb{R}}\Lambda^s( u^2)\Lambda^su \,\mathrm{d}x
-\lambda_2\int_{\mathbb{R}}\Lambda^s( u u_x)\Lambda^su \,\mathrm{d}x \\
&\quad-\int_{\mathbb{R}}\Lambda^s( G\ast(\lambda_3 u^2 + \lambda_4 u_x^2))\Lambda^su \,\mathrm{d}x
-\lambda_1\int_{\mathbb{R}}\Lambda^s( \partial_x \left(  G \ast (\lambda_5 u^2 + \lambda_6 u_x^2) \right))\Lambda^su \,\mathrm{d}x.
\end{aligned}
\end{equation}
Then, using the H\"{o}lder inequality, the Young inequality, the Kato-Ponce inequalities \eqref{KP1} and \eqref{KP2}, we estimate the four terms on the right hand side of \eqref{Hsblow-up} as follows:
\begin{equation}\label{Term1}
\begin{aligned}
&\quad\left|\lambda_1\int_{\mathbb{R}}\Lambda^s( u^2)\Lambda^su \,\mathrm{d}x\right|\leq c\|\Lambda^s(u^2)\|_{L^2}\|\Lambda^su\|_{L^2}\\
&\leq c\|u\|_{L^\infty}\|u\|_{H^s}^2,
\end{aligned}
\end{equation}
\begin{equation}\label{Term2}
\begin{aligned}
&\quad\left|\lambda_2\int_{\mathbb{R}}\Lambda^s( uu_x)\Lambda^su \,\mathrm{d}x\right|\\
&\leq\left|\lambda_2\int_{\mathbb{R}}\left(\Lambda^s( uu_x)- u \Lambda^s u_x\right)\Lambda^su \,\mathrm{d}x\right|+\left|\lambda_2\int_{\mathbb{R}} u \Lambda^s u_x\Lambda^su \,\mathrm{d}x\right|\\
&\leq c\|[\Lambda^s,u]u_x\|_{L^2} \|\Lambda^su\|_{L^2}+\left|\lambda_2\int_{\mathbb{R}} u_x (\Lambda^su)^2 \,\mathrm{d}x\right|\\
&\leq c\|\partial_x u\|_{L^\infty}\|u\|_{H^s}^2,
\end{aligned}
\end{equation}
\begin{equation}\label{Term3}
\begin{aligned}
&\quad\left|\int_{\mathbb{R}}\Lambda^s( G\ast(\lambda_3 u^2 + \lambda_4 u_x^2))\Lambda^su \,\mathrm{d}x\right|\\
&=\left|\int_{\mathbb{R}}\Lambda^{s-1}( \lambda_3 u^2 + \lambda_4 u_x^2)\Lambda^{s-1}u \,\mathrm{d}x\right|\\
&\leq c(\|u\|_{L^\infty}\|u\|_{H^{s-1}}+\|u_x\|_{L^\infty}\|u_x\|_{H^{s-1}})\|u\|_{H^{s-1}}\\
&\leq c(\|u\|_{L^\infty}+\|u_x\|_{L^\infty})\|u\|_{H^{s}}^2,
\end{aligned}
\end{equation}
and
\begin{equation}\label{Term4}
\begin{aligned}
&\quad\left|\int_{\mathbb{R}}\Lambda^s( \partial_x(G\ast(\lambda_5 u^2 + \lambda_6 u_x^2)))\Lambda^su \,\mathrm{d}x\right|\\
&=\left|\int_{\mathbb{R}}\Lambda^{s-1}( \lambda_5 u^2 + \lambda_6 u_x^2)\Lambda^{s}u \,\mathrm{d}x\right|\\
&\leq c(\|u\|_{L^\infty}\|u\|_{H^{s-1}}+\|u_x\|_{L^\infty}\|u_x\|_{H^{s-1}})\|u\|_{H^{s}}\\
&\leq c(\|u\|_{L^\infty}+\|u_x\|_{L^\infty})\|u\|_{H^{s}}^2.
\end{aligned}
\end{equation}
Combining \eqref{Term1}-\eqref{Term4} in \eqref{Hsblow-up}, we know that
\begin{equation*}
\frac{\mathrm{d}}{\mathrm{d}t}\|u\|_{H^s}^2\leq c(\|u\|_{L^\infty}+\|u_x\|_{L^\infty})\|u\|_{H^{s}}^2.
\end{equation*}
By Gronwall's inequality and the condition \eqref{blsc}, we get
\begin{equation*}
\|u\|_{H^s}^2\leq \mathrm{e}^{cMt}\|u_0\|_{H^s}^2, \quad t\in [0,T).
\end{equation*}
This completes the proof of Theorem \ref{bls1}.
\end{proof}

\end{document}
